\begin{acknowledgement}
临近截稿，已为毕业设计辛勤数周的我坐在电脑前，毕业论文仅剩致谢落笔成文，怀着激动而感慨的心情写下这段文字。回忆四年前，坐在良乡校区长椅上吹着秋风感慨的少年，是否能想到自己将会走上一条怎样曲折不凡的道路呢。

感谢付小雁老师。您凭借在位姿估计领域深厚的学术积淀，为我解惑释疑、点拨思路；更在治学态度与人生规划上悉心指引，在我对未来方向迷茫不定时为我拨开迷雾、认可我的科研潜力，让我跳出竞赛的局限，坚定踏上学术探索之路。

我应感谢WPIE学生创新实验室。它是我大学四年乃至人生中一个重要的转折点，大一时青涩懵懂又充满激情的我来到实验室中，历经了一系列培养和竞赛的磨练，获得了难以计量的知识和能力上的提升，更为我的人生带来一种从未设想的可能性。

感谢指导老师张盛博老师，若没有您经营的实验室、提供的良好科研环境，没有您维护实验室的纽带。我将鲜有机会接触如此丰富的、与自己专业不同方向的知识。也难以在RoboMaster这个竞赛机会中，以工程维度思考诸多问题，让自己的设计真正走进实际场景。

感谢黄恩浩学长，你的无私帮助让我一个电子信息工程专业的学生，对软件工程领域形成了丰富的了解，作为我实验室成长路上的引路人与前辈，你给予了我莫大的帮助与指引；感谢竞赛时和我一起攻坚克难的队友们，一个人可以走得很快，一群人才能走得更远。我们并肩作战，同心协力，带领战队走向历史上的最高辉煌。

感谢我的家人、挚友与挚爱。历尽千帆后回首望去，方知美好就在脚下。人生不仅仅只有事业和工作，所有共同经历的每一瞬都值得被纪念和感恩，那些年的喜悦与欢乐、悲伤和泪水、成功或不甘，共同铸就了我青春最难忘的回忆。祝我们永远年轻，永远热烈，永远仰望星空。
\end{acknowledgement}